%{{{ preamble
\documentclass[a4paper,12pt]{article}
\usepackage{anysize}
\marginsize{2cm}{2cm}{1cm}{1cm}
%0.7\textwidth 6.0in \textheight = 664pt
\usepackage{xltxtra}
\usepackage{xunicode}
\usepackage{graphicx}
\usepackage{color}
\usepackage{xgreek}
\usepackage{fancyvrb}
\usepackage{minted}
\usepackage{listings}
\usepackage{enumitem} 
\usepackage{framed} 
\usepackage{relsize}
\usepackage{float} 
\usepackage{pstricks}
\usepackage{pst-node}
\usepackage{pst-blur}
\setmainfont[Mapping=tex-text]{FreeSerif}
%}}}
\begin{document}

\def\thesection {\arabic{section}. }
\def\thesubsection {\Roman{subsection}) }

\begin{titlepage}
    \begin{center}
        \begin{figure}[h] 
            \includegraphics[width=0.2\textwidth]{title/ntua_logo}
        \end{figure}
        \vspace{1cm}
        \begin{LARGE}\textbf{ΕΘΝΙΚΟ ΜΕΤΣΟΒΙΟ ΠΟΛΥΤΕΧΝΕΙΟ\\[1.5cm]}\end{LARGE}
        \begin{Large}
            ΣΧΟΛΗ ΗΜ\&ΜΥ\\[2cm]
            Προηγμένα Θέματα\\
            Αρχιτεκτονικής Υπολογιστών\\[2cm]
            1\textsuperscript{η} Άσκηση\\
            Ακ. έτος 2011-2012\\
        \end{Large}
        \vfill
        \begin{flushright}
            \begin{tabular}{l r}
                {Γρηγόρης \textsc{Λύρας}}&
                {Α.Μ.: 03109687}\\
            \end{tabular}
        \end{flushright}

        \large\today\\
    \end{center}
\end{titlepage}




\section*{Εισαγωγή}

%{{{ Προσομοίωση

\subsection{Προσομοίωση}
Ο κώδικας που μας δόθηκε ήταν ο ακόλουθος:

\inputminted[linenos,fontsize=\scriptsize,frame=leftline]{c}{files/partA1-default-partA.c}

%}}}

\pagebreak

%{{{ Ιεραρχία μνήμης και μοντέλο απόδοσης

\subsection{Ιεραρχία μνήμης και μοντέλο απόδοσης}

Χρησιμοποιήσαμε την ιεραρχία μνήμης όπως μας δίνεται στο Παράρτημα.

\begin{table}[H]
    \centering
    \begin{tabular}{| c | c | c | c | c |}
        \hline
           & assoc  &   line size &   lines & size\\
        \hline
        \hline
        L1 instruction cache   & 2  &   64  &   512 & 32768 = 32 KB\\
        \hline
        L1 data cache   & 2  &   64  &   512 & 32768 = 32 KB\\
        \hline
        L2 cache  & 4  &   128  &   1024 & 131072 = 128 KB\\
        \hline
    \end{tabular}
    \caption{Cache Hierarchy}
    \label{fig:T2}
\end{table}


\begin{table}[H]
    \centering
    \begin{tabular}{| c | c |}
        \hline
        Cycles & 1885656236 \\
        \hline
        L1 miss ratio & 0.0117 \\
        \hline
        L2 miss ratio & 0.0183 \\
        \hline
    \end{tabular}
    \caption{Μετρικές πρώτης εκτέλεσης}
    \label{fig:T3}
\end{table}

%}}}

\section*{Τεχνικές Βελτιστοποίησης}
%{{{ Loop Interchange

\subsection{Loop Interchange}

Για την προσομοίωση κάναμε 6 εκτελέσεις για κάθε διάταξη των i,j,k όπως
φαίνεται στον πίνακα \ref{fig:T1}.

\begin{table}[H]
    \centering
    \begin{tabular}{| c | c |}
        \hline
        Σειρά προσομοίωσης & Διάταξη \\
        \hline
        \hline
        1   &  i,j,k \\
        \hline
        2   &  i,k,j \\
        \hline
        3   &  j,i,k \\
        \hline
        4   &  j,k,i \\
        \hline
        5   &  k,i,j \\
        \hline
        6   &  k,j,i \\
        \hline
    \end{tabular}
    \caption{Δυνατές Διατάξεις}
    \label{fig:T1}
\end{table}


Τα αποτελέσματα φαίνονται στα ακόλουθα σχήματα.

\begin{figure}[H]
	\centering
    \includegraphics[width=0.8\textwidth]{files/partA1-simics-outputs-cycles.png}
	\caption{Cycles}
	\label{fig:A1}
\end{figure}

\begin{figure}[H]
	\centering
    \includegraphics[width=0.8\textwidth]{files/partA1-simics-outputs-l1_miss.png}
    \caption{L1 Miss Rate}
	\label{fig:A2}
\end{figure}

\begin{figure}[H]
	\centering
    \includegraphics[width=0.8\textwidth]{files/partA1-simics-outputs-l2_miss.png}
    \caption{L2 Miss Rate}
	\label{fig:A3}
\end{figure}



%}}}


\pagebreak

\subsection{Cache Blocking}

Η L1 cache έχει μέγεθος 32KB, και line size 64 bytes. Κάνουμε πράξεις μεταξύ
αριθμών κινητής υποδιαστολής σε 32bit σύστημα συνεπώς κάθε ένας έχει μέγεθος
4bytes. Κάθε cache line χωράει 16 αριθμούς συνεπώς για να εκμεταλλευτούμε
καλύτερα την τοπικότητα των αναφορών θα εκτελούμε κάθε loop σε 16δες. Έτσι θα
χρησιμοποιούμε δεδομένα που έχουν ήδη έρθει στην cache.

\inputminted[linenos,fontsize=\scriptsize,frame=leftline]{c}{files/partA2-v3_patched-partA.c}


\begin{table}[H]
    \centering
    \begin{tabular}{| c | c |}
        \hline
        Cycles & 1448776923 \\
        Cycles & 9385880    \\
        \hline
        L1 miss ratio & 0.0131 \\
        L1 miss ratio & 0.0026 \\
        \hline
        L2 miss ratio & 0.0292 \\
        L2 miss ratio & 0.0065 \\
        \hline
    \end{tabular}
    \caption{Διαφορά των εκτελέσεων}
    \label{fig:T4}
\end{table}

\begin{figure}[H]
	\centering
    \includegraphics[width=0.8\textwidth]{files/partA2-simics-outputs-cycles.png}
	\caption{Cycles}
	\label{fig:B1}
\end{figure}

\begin{figure}[H]
	\centering
    \includegraphics[width=0.8\textwidth]{files/partA2-simics-outputs-l1_miss.png}
    \caption{L1 Miss Rate}
	\label{fig:B2}
\end{figure}

\begin{figure}[H]
	\centering
    \includegraphics[width=0.8\textwidth]{files/partA2-simics-outputs-l2_miss.png}
    \caption{L3 Miss Rate}
	\label{fig:B3}
\end{figure}



%}}}

\end{document}
